\section{From local contracts to embedded derivations}
\label{sec:internalization}

This section proves the core compiler theorem at the human level and matches
its pieces to the current Lean API.  The proof has two independent parts: a
recursive component theorem, already available in checked form in the
repository, and a finite topological composition theorem.

\subsection{Finite-trace reflection}

Let \(f:\N^k\to\N\) and \(g:\N^{k+2}\to\N\) be fixed primitive-recursive
functions.  Write
\[
 r(0,\vec x)=f(\vec x),\qquad
 r(t+1,\vec x)=g(t,r(t,\vec x),\vec x).
\]
The symbol for \(r\) is an instance of primitive recursion.

\begin{theorem}[Finite-trace reflection]
\label{thm:trace-reflection}
Let \(I(s,t,\vec x)\) be quantifier-free.  Suppose iPRA has embedded
derivations, under a context \(\Gamma(\vec x)\), of
\begin{align}
 I(f(\vec x),0,\vec x),&\label{eq:trace-base}\\
 I(s,t,\vec x)\longrightarrow
 I(g(t,s,\vec x),t+1,\vec x).&\label{eq:trace-step}
\end{align}
Then it has an embedded derivation of
\[
 \forall t\;I(r(t,\vec x),t,\vec x).
\]
For every PR fuel term \(F(\vec x)\), it consequently derives
\[
 I(r(F(\vec x),\vec x),F(\vec x),\vec x).
\]
\end{theorem}

\begin{proof}
Apply quantifier-free induction to
\(J(t,\vec x)\defeq I(r(t,\vec x),t,\vec x)\).  Substitution preserves the
quantifier-free property.  At zero, the defining equation for primitive
recursion rewrites \(r(0,\vec x)\) to \(f(\vec x)\), after which
\eqref{eq:trace-base} applies.  For the successor step, assume \(J(t,\vec x)\),
instantiate \eqref{eq:trace-step} with \(s=r(t,\vec x)\), and use the successor
defining equation to rewrite
\(r(t+1,\vec x)\) to \(g(t,r(t,\vec x),\vec x)\).  Universal elimination at
\(F(\vec x)\) gives the specialized statement.
\end{proof}

This proof is exactly the mathematical content of the repository's checked
\path{primitiveRecursionInvariant} and \path{primitiveRecursionAt}
constructors.  The important point is that the premises in a trace package are
object derivations.  The theorem does not inspect whether the invariant is
true in the standard model.

\subsection{Terminal extraction}

\begin{corollary}[Trace finalization]
\label{cor:trace-finalization}
Under the hypotheses of Theorem~\ref{thm:trace-reflection}, suppose iPRA also
derives
\[
 I(r(F(\vec x),\vec x),F(\vec x),\vec x)\longrightarrow B(\vec x).
\]
Then iPRA derives \(B(\vec x)\).
\end{corollary}

\begin{proof}
Specialize the universal invariant at \(F\) and apply implication elimination.
In the embedded calculus this is one use of simultaneous cut or two natural
deduction constructors.
\end{proof}

The checked stable-observables module already implements this pattern as
\path{loopProducesCertificate}.  The PCOG Lean core packages the same
operation as \path{TraceExit.compile}.

\subsection{Topological proof composition}

Let \(G\) be a finite topologically ordered graph with node outputs
\(B_0,\ldots,B_{N-1}\).  Node \(i\) names a list
\(R_i\subseteq\{B_0,\ldots,B_{i-1}\}\) and carries a proof constructor which,
for every context \(\Gamma\), maps proofs of all formulas in \(R_i\) to a
proof of \(B_i\).

\begin{theorem}[Finite proof-graph compilation]
\label{thm:proof-graph}
From the node constructors of a finite topologically ordered proof graph and
embedded proofs of its initial assumptions, one can construct an embedded
proof of every node output, in particular the designated final output.
\end{theorem}

\begin{proof}
Induct on the node list.  The empty prefix has no outputs.  Suppose proofs have
been constructed for all outputs in the prefix.  Every requirement of the
next node occurs in that prefix by the topological-order condition, so retrieve
its proof and feed the finite family to the node constructor.  Add the new
proof to the prefix table.  After \(N\) steps, the final formula has a proof.
This induction is performed in the host proof assistant over the finite graph;
it is not an induction axiom inside iPRA.  Each generated edge is an ordinary
composition of embedded \texttt{Prf} terms.
\end{proof}

The accompanying Lean source gives the intended implementation.  A
\texttt{Compiled} prefix stores a list of formulas and a function proving any
member.  \texttt{Compiled.extend} performs the inductive step.  The design
requires no decidable equality on formulas beyond the ordinary elimination of
list membership.

\subsection{Staged internalization theorem}

We can now state the main result.

\begin{theorem}[Staged internalization]
\label{thm:staged-internalization}
Let \(G\) be a fixed finite PCOG over free numerical parameters \(\vec x\).
Assume:
\begin{enumerate}
\item every function appearing in \(G\) has been compiled to a concrete
symbol of \texttt{PR};
\item every bounded predicate has a Boolean reflection package as in
Theorem~\ref{thm:bounded-booleanization};
\item every exact node carries an embedded proof of its local equations or
contract;
\item every trace node carries a trace package and a finalizer satisfying
Theorem~\ref{thm:trace-reflection};
\item every relational node carries embedded proofs of its finite validity,
compatibility, and bounded-search contracts;
\item every observation and finalizer node carries the stated embedded local
proof constructor;
\item all dependency and demand edges respect the graph's topological and
resource contracts.
\end{enumerate}
Then there is an effective host construction of a term
\[
  \Prf\;\PRAAx\;\Gamma\;\Phi_G,
\]
where \(\Phi_G\) is the final graph formula.  If \(\Gamma\) consists of input
validity assumptions, the assumptions can be universally closed and introduced
as implications to obtain a closed theorem.
\end{theorem}

\begin{proof}
First compile every trace node independently using
Theorem~\ref{thm:trace-reflection} and its finalizer.  This turns each recursive
component into an ordinary proof-producing rule.  Exact, relational, and
observation nodes already have that form.  The component graph is finite and
acyclic by construction, so apply Theorem~\ref{thm:proof-graph}.  The result is
an embedded derivation under \(\Gamma\).  Repeated universal introduction and
implication introduction close the free parameters and hypotheses, using the
scoped substitution/weakening lemmas of the backend.
\end{proof}

The theorem is ``effective'' in a precise metatheoretic sense: a Lean function
can build the proof term from finite node packages.  It is not asserting that
an iPRA program can interpret an arbitrary encoded graph.

\subsection{Why no universal evaluator is needed}

A total primitive-recursive universal evaluator for all primitive-recursive
program codes cannot exist.  If \(U(e,x)\) were PR and represented every unary
PR function, then
\[
 d(x)=U(x,x)+1
\]
would be PR and would differ from the function represented by its own code at
that code.  This diagonal contradiction is elementary and formalizable at the
metalevel.

The staged compiler avoids the premise.  Given a fixed node syntax tree, the
host recursively constructs a \texttt{PR} value by the inductive constructors
\texttt{zero}, \texttt{succ}, \texttt{proj}, \texttt{comp}, and
\texttt{prec}.  The resulting arithmetic theorem mentions that one concrete
symbol.  Another graph produces another symbol.  There is no \(U\) taking a
runtime code.

\begin{principle}[Specialize before internalizing]
Whenever a proposed framework seems to require interpreting arbitrary program
syntax inside iPRA, move the interpretation to the host compiler and emit
concrete \texttt{PR} syntax.  Quantification over numerical \emph{inputs} is
allowed; quantification over arbitrary program \emph{codes} is not needed.
\end{principle}

\subsection{Composition and substitution}

Suppose one compiled graph proves \(A(\vec x)\), and a second graph, under
assumption \(A(\vec y)\), proves \(B(\vec y)\).  A term substitution
\(\vec y=\vec t(\vec x)\) turns the second derivation into one under
\(A(\vec t(\vec x))\).  If the first graph exports that instance, simultaneous
cut yields \(B(\vec t(\vec x))\).  This is the formal content of contract
composition.

Precision composition occurs separately: the second graph's backward pass
computes demands on \(A\)'s observations, which become output requests to the
first graph.  Keeping logical composition and precision composition separate
is valuable.  The former is proof substitution and cut; the latter is a
primitive-recursive numerical program whose soundness is one of the local
contracts.

\subsection{Acyclic normalization of complex algorithms}

An arbitrary fixed algorithm built from primitive recursion, bounded search,
and finite composition can be converted to PCOG normal form:
\begin{enumerate}
\item identify maximal uniform recurrences and package each as a trace node;
\item replace bounded existential choices by least-search or relational cloud
nodes;
\item isolate all error estimates and algebraic identities as local proof
nodes;
\item collapse witness-dependent outputs to stable observations only after a
comparison theorem;
\item connect the resulting components by an acyclic dependency graph.
\end{enumerate}

Nested loops do not require nested logical induction if their inner run is
compiled to a PR symbol with a proved contract; the outer invariant invokes
that contract as a node rule.  When the inner loop's correctness depends on an
outer state, its theorem is uniformly parameterized by that state.

This normal form is the proposed organizing principle for future theorem
packets.  It gives a standard answer to ``where is the actual iPRA proof?'':
it is the result of compiling the trace packages and folding the proof graph.
